% Written by the tex module's build_pdf from references.toml: edit the records, not this file.
\providecommand{\url}[1]{\texttt{#1}}\providecommand{\nolinkurl}[1]{\texttt{#1}}
\begin{thebibliography}{46}
\bibitem{lefebvre1995} Lefebvre S, Bürglen L, Reboullet S, Clermont O, Burlet P, Viollet L, et al. Identification and characterization of a spinal muscular atrophy-determining gene. Cell 1995;80(1):155-65. doi:\nolinkurl{10.1016/0092-8674(95)90460-3}
\bibitem{lorson1999} Lorson CL, Hahnen E, Androphy EJ, Wirth B. A single nucleotide in the SMN gene regulates splicing and is responsible for spinal muscular atrophy. Proc Natl Acad Sci U S A 1999;96(11):6307-11. doi:\nolinkurl{10.1073/pnas.96.11.6307}
\bibitem{monani1999} Monani UR, Lorson CL, Parsons DW, Prior TW, Androphy EJ, Burghes AH, et al. A single nucleotide difference that alters splicing patterns distinguishes the SMA gene SMN1 from the copy gene SMN2. Hum Mol Genet 1999;8(7):1177-83. doi:\nolinkurl{10.1093/hmg/8.7.1177}
\bibitem{mercuri2022} Mercuri E, Sumner CJ, Muntoni F, Darras BT, Finkel RS. Spinal muscular atrophy. Nat Rev Dis Primers 2022;8(1):52. doi:\nolinkurl{10.1038/s41572-022-00380-8}
\bibitem{finkel2017} Finkel RS, Mercuri E, Darras BT, Connolly AM, Kuntz NL, Kirschner J, et al. Nusinersen versus Sham Control in Infantile-Onset Spinal Muscular Atrophy. N Engl J Med 2017;377(18):1723-32. doi:\nolinkurl{10.1056/NEJMoa1702752}
\bibitem{mendell2017} Mendell JR, Al-Zaidy S, Shell R, Arnold WD, Rodino-Klapac LR, Prior TW, et al. Single-Dose Gene-Replacement Therapy for Spinal Muscular Atrophy. N Engl J Med 2017;377(18):1713-22. doi:\nolinkurl{10.1056/NEJMoa1706198}
\bibitem{prior2008} Prior TW. Carrier screening for spinal muscular atrophy. Genet Med 2008;10(11):840-2. doi:\nolinkurl{10.1097/GIM.0b013e318188d069}
\bibitem{hendrickson2009} Hendrickson BC, Donohoe C, Akmaev VR, Sugarman EA, Labrousse P, Boguslavskiy L, et al. Differences in SMN1 allele frequencies among ethnic groups within North America. J Med Genet 2009;46(9):641-4. doi:\nolinkurl{10.1136/jmg.2009.066969}
\bibitem{sugarman2012} Sugarman EA, Nagan N, Zhu H, Akmaev VR, Zhou Z, Rohlfs EM, et al. Pan-ethnic carrier screening and prenatal diagnosis for spinal muscular atrophy: clinical laboratory analysis of >72,400 specimens. Eur J Hum Genet 2012;20(1):27-32. doi:\nolinkurl{10.1038/ejhg.2011.134}
\bibitem{luo2014} Luo M, Liu L, Peter I, Zhu J, Scott SA, Zhao G, et al. An Ashkenazi Jewish SMN1 haplotype specific to duplication alleles improves pan-ethnic carrier screening for spinal muscular atrophy. Genet Med 2014;16(2):149-56. doi:\nolinkurl{10.1038/gim.2013.84}
\bibitem{feng2017} Feng Y, Ge X, Meng L, Scull J, Li J, Tian X, et al. The next generation of population-based spinal muscular atrophy carrier screening: comprehensive pan-ethnic SMN1 copy-number and sequence variant analysis by massively parallel sequencing. Genet Med 2017;19(8):936-44. doi:\nolinkurl{10.1038/gim.2016.215}
\bibitem{milligan2021} Milligan JN, Larson JL, Filipovic-Sadic S, Laosinchai-Wolf W, Huang YW, Ko TM, et al. Multisite Evaluation and Validation of a Sensitive Diagnostic and Screening System for Spinal Muscular Atrophy that Reports SMN1 and SMN2 Copy Number, along with Disease Modifier and Gene Duplication Variants. J Mol Diagn 2021;23(6):753-64. doi:\nolinkurl{10.1016/j.jmoldx.2021.03.004}
\bibitem{mikhalchuk2025} Mikhalchuk K, Polyakov A, Zabnenkova VV, Ismagilova O, Shchagina O. Exploring the Prevalence of SMN1 Duplication and Deletion in Russia and Its Impact on Carrier Screening. Int J Mol Sci 2025;26(5). doi:\nolinkurl{10.3390/ijms26051984}
\bibitem{chen2020} Chen X, Sanchis-Juan A, French CE, Connell AJ, Delon I, Kingsbury Z, et al. Spinal muscular atrophy diagnosis and carrier screening from genome sequencing data. Genet Med 2020;22(5):945-53. doi:\nolinkurl{10.1038/s41436-020-0754-0}
\bibitem{chen2023} Chen X, Harting J, Farrow E, Thiffault I, Kasperaviciute D, Hoischen A, et al. Comprehensive SMN1 and SMN2 profiling for spinal muscular atrophy analysis using long-read PacBio HiFi sequencing. Am J Hum Genet 2023;110(2):240-50. doi:\nolinkurl{10.1016/j.ajhg.2023.01.001}
\bibitem{brzustowicz1990} Brzustowicz LM, Lehner T, Castilla LH, Penchaszadeh GK, Wilhelmsen KC, Daniels R, et al. Genetic mapping of chronic childhood-onset spinal muscular atrophy to chromosome 5q11.2-13.3. Nature 1990;344(6266):540-1. doi:\nolinkurl{10.1038/344540a0}
\bibitem{melki1990} Melki J, Abdelhak S, Sheth P, Bachelot MF, Burlet P, Marcadet A, et al. Gene for chronic proximal spinal muscular atrophies maps to chromosome 5q. Nature 1990;344(6268):767-8. doi:\nolinkurl{10.1038/344767a0}
\bibitem{gilliam1990} Gilliam TC, Brzustowicz LM, Castilla LH, Lehner T, Penchaszadeh GK, Daniels RJ, et al. Genetic homogeneity between acute and chronic forms of spinal muscular atrophy. Nature 1990;345(6278):823-5. doi:\nolinkurl{10.1038/345823a0}
\bibitem{roy1995} Roy N, Mahadevan MS, McLean M, Shutler G, Yaraghi Z, Farahani R, et al. The gene for neuronal apoptosis inhibitory protein is partially deleted in individuals with spinal muscular atrophy. Cell 1995;80(1):167-78. doi:\nolinkurl{10.1016/0092-8674(95)90461-1}
\bibitem{burglen1997} Bürglen L, Seroz T, Miniou P, Lefebvre S, Burlet P, Munnich A, et al. The gene encoding p44, a subunit of the transcription factor TFIIH, is involved in large-scale deletions associated with Werdnig-Hoffmann disease. Am J Hum Genet 1997;60(1):72-9.
\bibitem{campbell1997} Campbell L, Potter A, Ignatius J, Dubowitz V, Davies K. Genomic variation and gene conversion in spinal muscular atrophy: implications for disease process and clinical phenotype. Am J Hum Genet 1997;61(1):40-50. doi:\nolinkurl{10.1086/513886}
\bibitem{vandersteege1996} van der Steege G, Grootscholten PM, Cobben JM, Zappata S, Scheffer H, den Dunnen JT, et al. Apparent gene conversions involving the SMN gene in the region of the spinal muscular atrophy locus on chromosome 5. Am J Hum Genet 1996;59(4):834-8.
\bibitem{vorster2020} Vorster E, Essop FB, Rodda JL, Krause A. Spinal Muscular Atrophy in the Black South African Population: A Matter of Rearrangement?. Front Genet 2020;11:54. doi:\nolinkurl{10.3389/fgene.2020.00054}
\bibitem{wadman2020} Wadman RI, Jansen MD, Stam M, Wijngaarde CA, Curial CAD, Medic J, et al. Intragenic and structural variation in the SMN locus and clinical variability in spinal muscular atrophy. Brain Commun 2020;2(2):fcaa075. doi:\nolinkurl{10.1093/braincomms/fcaa075}
\bibitem{berzal2026} Berzal-Serrano A, González-Alvaredo S, Lesende I, Carretero-Vilarroig L, Jaijo T, Onsurbe I, et al. Multimodal characterisation of the SMN locus in SMA: copy number quantification and hybrid gene identification. NPJ Genom Med 2026;11(1). doi:\nolinkurl{10.1038/s41525-026-00579-8}
\bibitem{liao2023} Liao WW, Asri M, Ebler J, Doerr D, Haukness M, Hickey G, et al. A draft human pangenome reference. Nature 2023;617(7960):312-24. doi:\nolinkurl{10.1038/s41586-023-05896-x}
\bibitem{lucas2026} Lucas JK, Hebbar P, Liao W, Macias-Velasco JF, Novak AM, Asri M, et al. HPRC2: A human pangenome reference with near-complete coverage of common genetic variation. bioRxiv (preprint) 2026. doi:\nolinkurl{10.64898/2026.07.21.739710}
\bibitem{logsdon2025} Logsdon GA, Ebert P, Audano PA, Loftus M, Porubsky D, Ebler J, et al. Complex genetic variation in nearly complete human genomes. Nature 2025;644(8076):430-41. doi:\nolinkurl{10.1038/s41586-025-09140-6}
\bibitem{suzuki2026} Suzuki Y, Owa C, Kobayashi H, Nakabayashi R, McNulty B, Violich I, et al. Accessing medically relevant complex regions with a pangenome graph of 20 near-complete Japanese haplotypes. Nat Commun 2026;17(1). doi:\nolinkurl{10.1038/s41467-026-73461-x}
\bibitem{zwartkruis2025} Zwartkruis MM, Elferink MG, Gommers D, Signoria I, Blasco-Pérez L, Costa-Roger M, et al. Long-read sequencing identifies copy-specific markers of SMN gene conversion in spinal muscular atrophy. Genome Med 2025;17(1):26. doi:\nolinkurl{10.1186/s13073-025-01448-2}
\bibitem{hickey2024} Hickey G, Monlong J, Ebler J, Novak AM, Eizenga JM, Gao Y, et al. Pangenome graph construction from genome alignments with Minigraph-Cactus. Nat Biotechnol 2024;42(4):663-73. doi:\nolinkurl{10.1038/s41587-023-01793-w}
\bibitem{chin2023} Chin CS, Behera S, Khalak A, Sedlazeck FJ, Sudmant PH, Wagner J, et al. Multiscale analysis of pangenomes enables improved representation of genomic diversity for repetitive and clinically relevant genes. Nat Methods 2023;20(8):1213-21. doi:\nolinkurl{10.1038/s41592-023-01914-y}
\bibitem{anderson2001} Anderson MJ. A new method for non-parametric multivariate analysis of variance. Austral Ecology 2001;26:32-46. doi:\nolinkurl{10.1111/j.1442-9993.2001.01070.pp.x}
\bibitem{yoo2025} Yoo D, Rhie A, Hebbar P, Antonacci F, Logsdon GA, Solar SJ, et al. Complete sequencing of ape genomes. Nature 2025;641(8062):401-18. doi:\nolinkurl{10.1038/s41586-025-08816-3}
\bibitem{rochette2001} Rochette CF, Gilbert N, Simard LR. SMN gene duplication and the emergence of the SMN2 gene occurred in distinct hominids: SMN2 is unique to Homo sapiens. Hum Genet 2001;108(3):255-66. doi:\nolinkurl{10.1007/s004390100473}
\bibitem{rozen2003} Rozen S, Skaletsky H, Marszalek JD, Minx PJ, Cordum HS, Waterston RH, et al. Abundant gene conversion between arms of palindromes in human and ape Y chromosomes. Nature 2003;423(6942):873-6. doi:\nolinkurl{10.1038/nature01723}
\bibitem{bonito2021} Bonito M, D'Atanasio E, Ravasini F, Cariati S, Finocchio A, Novelletto A, et al. New insights into the evolution of human Y chromosome palindromes through mutation and gene conversion. Hum Mol Genet 2021;30(23):2272-85. doi:\nolinkurl{10.1093/hmg/ddab189}
\bibitem{bolognini2024} Bolognini D, Halgren A, Lou RN, Raveane A, Rocha JL, Guarracino A, et al. Recurrent evolution and selection shape structural diversity at the amylase locus. Nature 2024;634(8034):617-25. doi:\nolinkurl{10.1038/s41586-024-07911-1}
\bibitem{yilmaz2024} Yilmaz F, Karageorgiou C, Kim K, Pajic P, Scheer K, Beck CR, et al. Reconstruction of the human amylase locus reveals ancient duplications seeding modern-day variation. Science 2024;386(6724):eadn0609. doi:\nolinkurl{10.1126/science.adn0609}
\bibitem{katoh2013} Katoh K, Standley DM. MAFFT multiple sequence alignment software version 7: improvements in performance and usability. Mol Biol Evol 2013;30(4):772-80. doi:\nolinkurl{10.1093/molbev/mst010}
\bibitem{minh2020} Minh BQ, Schmidt HA, Chernomor O, Schrempf D, Woodhams MD, von Haeseler A, et al. IQ-TREE 2: New Models and Efficient Methods for Phylogenetic Inference in the Genomic Era. Mol Biol Evol 2020;37(5):1530-4. doi:\nolinkurl{10.1093/molbev/msaa015}
\bibitem{hoang2018} Hoang DT, Chernomor O, von Haeseler A, Minh BQ, Vinh LS. UFBoot2: Improving the Ultrafast Bootstrap Approximation. Mol Biol Evol 2018;35(2):518-22. doi:\nolinkurl{10.1093/molbev/msx281}
\bibitem{mantel1967} Mantel N. The detection of disease clustering and a generalized regression approach. Cancer Res 1967;27(2):209-20.
\bibitem{harris2020} Harris CR, Millman KJ, van der Walt SJ, Gommers R, Virtanen P, Cournapeau D, et al. Array programming with NumPy. Nature 2020;585(7825):357-62. doi:\nolinkurl{10.1038/s41586-020-2649-2}
\bibitem{virtanen2020} Virtanen P, Gommers R, Oliphant TE, Haberland M, Reddy T, Cournapeau D, et al. SciPy 1.0: fundamental algorithms for scientific computing in Python. Nat Methods 2020;17(3):261-72. doi:\nolinkurl{10.1038/s41592-019-0686-2}
\bibitem{matplotlib} The Matplotlib Development Team. Matplotlib: visualization with Python [computer program]. 2025 [cited 2026-11-14]. Available from: \url{https://matplotlib.org}
\end{thebibliography}
